\scanpage{41}
\mspage{40}
\mspage{81}
\margin{n°382}

$f\in J^k(X,K^q)_0$).

7.1.7. Si Y = K, la structure d'algèbre des $P_x^k(X)$ munit $P^k(X)$ d'une structure d'espace fibré vectoriel en algèbres de base X, de classe $C^{r-k}$.

7.1.8. Soit X une variété de classe $C^r$, et soit $x\in X$. L'espace vectoriel $J_0^1(K,X)_x$ s'identifie canoniquement à l'espace tangent $T_x(X)$ à X en x. La réunion $T(X)=\bigcup_{x\in X}T_x(X)$ des espaces tangents aux points de X s'identifie canoniquement à $J^1(K,X)$, et admet donc une structure d'espace fibré vectoriel sur X, de classe $C^{r-1}$. C'est le \uline{fibré tangent} $T(X)$ à X.

7.1.9. Soient X et Y des variétés de classe $C^r$ et k un entier fini $\leq r$. Un morphisme $f:X\longrightarrow Y$ définit une application $x\longrightarrow j_x^k(f)$ de X dans $J^k(X,Y)$ qui est une section de classe $C^{r-k}$ pour la fibration de $J^k(X,Y)$ sur X. L'application $J^k(X,Y)\longrightarrow J^{k'}(X,Y)$ (k' entier $\leq k$, cf. n° 7.1.2) est un morphisme pour les fibrations sur X et pour les fibrations sur Y.

7.1.10. Soient X, Y des variétés de classe $C^r$, $k'\leq k$ des entiers finis $\leq r$, $\ell=k-k'$. Soit f une application de classe $C^r$, dans Y, d'un ouvert U de X. L'application $x\longrightarrow j^{k'}(f)$ de U dans $J^{k'}(X,Y)$ est de classe $C^{r-k'}$ et son jet d'ordre $\ell$ en x ne dépend que de $j^k(f)$. On obtient ainsi une application canonique de classe $C^{r-k}$ de $J^k(X,Y)$ dans $J^\ell(X,J^{k'}(X,Y))$ qui, si K est de caractéristique zéro, est un isomorphisme de $J^k(X,Y)$ sur une sous-variété de $J^\ell(X,J^{k'}(X,Y))$ et qui, si X et Y sont des variétés pures, est un morphisme pour les fibrations sur X.

7.1.11. Soient X, Y, X', Y' des variétés de classe $C^r$, $f:X\longrightarrow X'$ et $g:Y'\longrightarrow Y$ des morphismes, $x\in X$, $y'\in Y'$, $x'=f(x)$, $y=g(y')$. Les morphismes f et g définissent une application $j_{x'}^k(f,g)_y:J_{x'}^k(X',Y')_{y'}\longrightarrow J_x^k(X,Y)_y$ caractérisée par $u\longrightarrow P_{y'}^k\circ u\circ P_x^k(f)$. On en déduit une application, $j^k(f,g)$ de $f^{-1}J^k(X',Y')=X\times_{X'}J^k(X',Y')$ dans $J^k(X,Y)$, de classe $C^{r-k}$. Si ces
